\subsubsection{Sparse Diffusion Targeting}

A key observation in our early experiments is that directly applying diffusion enhancement to the entire hard mask leads to severe overuse of generative modification. This overuse damages surrounding background consistency and introduces profound seam artifacts. Therefore, we explicitly restrict diffusion to a sparse subset of difficult pixels instead of the whole candidate region.

Our pipeline first applies ProPainter as the main restoration module, and then constructs a \texttt{diffusion\_target} mask. The target mask is obtained by subtracting the borrowable areas from the candidate regions, followed by confidence gating, morphological filtering, and connected-component cleanup. When the target region becomes too large, we enforce an upper bound on the ratio between diffusion-target pixels and hard-mask pixels (e.g., $0.35$ or $0.25$), retaining only the most difficult pixels. The difficulty score combines a low support fraction with high color inconsistency, prioritizing pixels that ProPainter failed to restore reliably.

This sparse design is strongly supported by our quantitative tracking. Here, \texttt{target/hard} denotes the ratio of pixels selected for diffusion refinement versus the total hard mask pixels. Before introducing the ratio constraint, the average \texttt{target/hard} ratio was $0.915$ (with a p90 of $0.994$), proving that diffusion was over-writing almost the entire hard region. After the sparse strategy was applied, we successfully clamped this ratio to $0.35$ and $0.25$ respectively.

\begin{table}[h]
\centering
\caption{Effect of sparse targeting on diffusion coverage. Lower values indicate more selective diffusion refinement.}
\label{tab:sparse_target}
\begin{tabular}{lcc}
\toprule
Setting & target/hard mean & target/hard p90 \\
\midrule
Before sparsification (run 130347) & 0.915 & 0.994 \\
Sparse target $0.35$ (run 133713) & 0.350 & 0.350 \\
Sparse target $0.25$ (run 135006) & 0.250 & 0.250 \\
\bottomrule
\end{tabular}
\end{table}

Among these settings, the sparse-$0.35$ configuration achieved the best balance between suppressing the over-generation risk and preserving keyframe stability, and was therefore established as our default setting.

\begin{figure}[t]
    \centering
    \includegraphics[width=\linewidth]{figures/part3_sparse_target_vis.png}
    \caption{Sparse diffusion targeting. From left to right: hard mask, diffusion target, and enhanced result.}
    \label{fig:sparse_target}
\end{figure}

\subsubsection{Improved Boundary Blending}

Besides the overuse of diffusion, another critical artifact natively introduced by the SDXL VAE decoder is the visible white "halo" (or bright boundary) around refined regions. Our analysis confirms that this issue was highly amplified by the rigid binary morphological mask (erosion/dilation) originally used for blending, which bluntly pasted the VAE padding artifacts back into the restored frame.

To eradicate this problem, we replaced the morphological approach with a Distance-Transform-based continuous alpha blending strategy. Instead of a hard transition, the new design calculates the exact Euclidean distance from the mask boundary and applies a perfect linear gradient ($\alpha \in [0,1]$). It enforces a $0.0$ alpha weight at the absolute boundary (preserving the pristine ProPainter background) and linearly climbs to $1.0$ towards the core.

In practice, this modification mathematically eradicated the halo artifacts, achieving a remarkable spatial improvement. As tracked in our metrics (Run 112532 vs Run 143539 generation), the outside change L1 mean ($\mathcal{L}_{1}^{\text{out}}$) plummeted by roughly 74\% (down to $0.0023$), and the seam change L1 mean ($\mathcal{L}_{1}^{\text{seam}}$) decreased by 56\% (down to $0.2841$). 

\begin{figure}[t]
    \centering
    \includegraphics[width=\linewidth]{results/part3/wildvideo/blending_comparison_labeled_horizontal.png}
    \caption{Improved boundary blending. Left: Original morphological hard-cut showing bright halos. Right: Distance-Transform alpha blending seamlessly hiding the VAE boundaries.}
    \label{fig:boundary_blending}
\end{figure}

\subsubsection{Failure Analysis of Flow-based Temporal Propagation}

Although localized diffusion dramatically improved spatial quality on individual keyframes, it introduced immense temporal inconsistency ($>3.72$ instability ratio) as the ultra-crisp generative textures alternated with softer ProPainter background frames in time. To mitigate this flickering, we aggressively explored optical-flow-based residual propagation, aiming to warp the high-definition textures from keyframes to neighboring non-keyframes.

However, cross-dataset evaluations ultimately proved this to be a theoretical dead end for complex real-world videos. Object removal inherently involves massive 3D parallax, non-rigid dynamic occlusions (e.g., foliage), and disocclusions. While 2D optical flow is adequate for tracking visible, rigid translations, it is mathematically incapable of hallucinating missing background depths or safely warping patches across severe occlusion boundaries. 

Our empirical data firmly corroborates this structural limitation.

\begin{table*}[t]
\centering
\small
\setlength{\tabcolsep}{4pt}
\caption{Cross-Dataset Evaluation of Residual Propagation via Optical Flow.}
\label{tab:flow_failure}
\begin{tabular}{llccc}
\toprule
Dataset & Flow Method & Hard Gain $\uparrow$ & Temporal Ratio $\downarrow$ & Seam $L_1 \downarrow$ \\
\midrule
\multirow{3}{*}{\texttt{wildvideo}} 
 & SD Base (No Prop) & 0.0419 & 3.720 & 0.6508 \\
 & + Farneback & 0.0695 & 3.724 & 0.6402 \\
 & + RAFT & 0.0000 & 3.907 & 0.3777 \\
\midrule
\multirow{2}{*}{\texttt{bmx-trees}} 
 & + Farneback & 0.0000 & 0.998 & 1.2241 \\
 & + RAFT & 0.0000 & 0.984 & 1.1177 \\
\bottomrule
\end{tabular}
\end{table*}

As shown in Table \ref{tab:flow_failure}, on the simpler \texttt{wildvideo} sequence, traditional Farneback flow managed to distribute textures and spike the Hard Gain to $0.0695$; however, temporal flickering remained severely high at $3.724$ due to sub-pixel misalignment accumulating over time, inducing texture tearing.

Crucially, on the complex \texttt{bmx-trees} dataset, the Farneback estimator failed completely (Gain $= 0.0$). The complex foliage and severe camera parallax induced completely invalid 2D warps, triggering the system's overlap validation threshold and discarding the frames wholesale. Furthermore, upgrading to the deep learning RAFT flow yielded universal gain annihilation ($0.0$) across all datasets. This exposed strict engineering vulnerabilities in cross-framework optical flow integrations, heavily penalizing mismatched tensor normalization ([-1, 1] vs [0, 255]) and reversed spatial grid coordinate mapping (Forward vs. Backward flow).

Overall, our experiments indicate that while the spatial domain operations (Sparse Targeting + Distance Alpha) masterfully resolved single-frame diffusion integration, temporal stability remains an insurmountable obstacle for post-generation 2D tracking. Future improvements must therefore discard post-hoc flow and embrace Joint Temporal Generation architectures directly utilizing 3D-UNets within the latent space.
